\section{第 12 章综合习题}

\begin{problem}{三角级数的一致收敛性与 $L^2$ 性质}{Trigonometric-Series-Uniform-Convergence}
    证明：级数 $\sum_{n=2}^{\infty} \frac{\sin nx}{\ln n}$ 在不包含 $2\uppi$ 整数倍的闭区间上一致收敛，但它不是任意在 $[-\uppi, \uppi]$ 上平方可积函数的 Fourier 级数．
\end{problem}

\begin{myproof}
    设 $[a, b]$ 是一个不包含 $2\uppi$ 整数倍的闭区间．此时，对任意 $x \in [a, b]$ 均有 $x \neq 2k\uppi$ ($k \in \symbb{Z}$)，从而 $\sin\frac{x}{2} \neq 0$．由连续函数的性质，存在正常数 $M$ 使得对任意 $x \in [a, b]$，均有 $\left|\sin\frac{x}{2}\right| \geqslant M > 0$．

    考虑三角部分的部和：
    \begin{equation}\label{eq:trig-bound}
        \left| \sum_{k=2}^{n} \sin kx \right| = \left| \frac{\cos\frac{3}{2}x - \cos\left(n+\frac{1}{2}\right)x}{2\sin\frac{x}{2}} \right| \leqslant \frac{1}{\left|\sin\frac{x}{2}\right|} \leqslant \frac{1}{M}.
    \end{equation}
    这说明 $\sum_{n=2}^{\infty} \sin nx$ 的部分和在区间 $[a, b]$ 上一致有界．又因为数列 $a_n = \frac{1}{\ln n}$ 单调递减且趋于 $0$，根据 Dirichlet 一致收敛判别法，级数 $\sum_{n=2}^{\infty} \frac{\sin nx}{\ln n}$ 在 $[a, b]$ 上一致收敛．

    接下来，假设该级数是某个在 $[-\uppi, \uppi]$ 上平方可积函数 $f(x) \in L^2([-\uppi, \uppi])$ 的 Fourier 级数，则其 Fourier 系数应为 $a_n = 0$ ($n \geqslant 0$)，$b_1 = 0$ 且 $b_n = \frac{1}{\ln n}$ ($n \geqslant 2$)．

    根据平方可积函数的 Bessel 不等式，其 Fourier 系数必须满足
    \begin{equation}\label{eq:bessel-ineq-proof}
        \sum_{n=2}^{\infty} b_n^2 = \sum_{n=2}^{\infty} \frac{1}{(\ln n)^2} \leqslant \frac{1}{\uppi} \int_{-\uppi}^{\uppi} |f(x)|^2 \d x < \infty.
    \end{equation}
    然而，当 $n \geqslant 2$ 时，显然有 $\ln n < n$，从而 $(\ln n)^2 < n^2$．实际上，由极限性质
    \begin{equation}\label{eq:limit-comp}
        \lim_{n \to \infty} \frac{\frac{1}{(\ln n)^2}}{\frac{1}{n}} = \lim_{n \to \infty} \frac{n}{(\ln n)^2} = \infty,
    \end{equation}
    可知对足够大的 $n$，有 $\frac{1}{(\ln n)^2} > \frac{1}{n}$．因为调和级数 $\sum_{n=2}^{\infty} \frac{1}{n}$ 发散，所以正项级数 $\sum_{n=2}^{\infty} \frac{1}{(\ln n)^2}$ 亦发散，这与\zcref{eq:bessel-ineq-proof}矛盾．

    因此，该级数不是任何在 $[-\uppi, \uppi]$ 上平方可积函数的 Fourier 级数．
\end{myproof}

\begin{problem}{三角级数求和恒等式}{Trigonometric-Series-Summation-Identities}
    证明下列等式：
    \begin{enumerate}
        \item $\sum_{n=1}^{\infty} (-1)^{n-1} \frac{\cos nx}{n} = \ln \left( 2\cos \frac{x}{2} \right) \quad (-\uppi < x < \uppi)$;
        \item $\sum_{n=1}^{\infty} \frac{\cos nx}{n} = -\ln \left( 2\sin \frac{x}{2} \right) \quad (0 < x < 2\uppi)$.
    \end{enumerate}
\end{problem}

\begin{myproof}
    \ansitem{2}

    设函数 $g(x) = -\ln \left( 2\sin \frac{x}{2} \right)$，其定义在区间 $(0, 2\uppi)$ 上．我们将其在 $[-\uppi, \uppi]$ 上开拓为偶函数，即令 $f(x) = -\ln \left( 2\left|\sin \frac{x}{2}\right| \right)$．
    由于 $f(x)$ 是偶函数，其 Fourier 级数中正弦系数 $b_n = 0$．下面计算余弦系数．
    首先计算 $a_0$：
    \begin{equation}\label{eq:a0-calc}
        a_0 = \frac{2}{\uppi} \int_{0}^{\uppi} -\ln \left( 2\sin \frac{x}{2} \right) \d x.
    \end{equation}
    令 $u = \frac{x}{2}$，则 $\d x = 2 \d u$，于是有
    \begin{equation}\label{eq:a0-integral}
        \int_{0}^{\uppi} \ln \left( 2\sin \frac{x}{2} \right) \d x = 2 \int_{0}^{\uppi/2} \left[ \ln 2 + \ln(\sin u) \right] \d u = \uppi \ln 2 + 2 \int_{0}^{\uppi/2} \ln(\sin u) \d u.
    \end{equation}
    利用已知积分 $\int_{0}^{\uppi/2} \ln(\sin u) \d u = -\frac{\uppi}{2} \ln 2$，可得
    \begin{equation}\label{eq:a0-zero}
        \int_{0}^{\uppi} \ln \left( 2\sin \frac{x}{2} \right) \d x = \uppi \ln 2 + 2 \left( -\frac{\uppi}{2} \ln 2 \right) = 0,
    \end{equation}
    因此 $a_0 = 0$．

    接下来对 $n \geqslant 1$，计算 $a_n$：
    \begin{equation}\label{eq:an-integral}
        a_n = \frac{2}{\uppi} \int_{0}^{\uppi} -\ln \left( 2\sin \frac{x}{2} \right) \cos(nx) \d x.
    \end{equation}
    应用分部积分法，令 $u = \ln \left( 2\sin \frac{x}{2} \right)$，$\d v = \cos(nx) \d x$，则 $\d u = \frac{1}{2} \cot \frac{x}{2} \d x$，$v = \frac{\sin(nx)}{n}$．
    由于在端点处有极限值 $\lim_{x \to 0^+} \ln \left( 2\sin \frac{x}{2} \right) \sin(nx) = 0$，故分部积分的第一项为零，从而有
    \begin{equation}\label{eq:an-parts}
        a_n = \frac{1}{n\uppi} \int_{0}^{\uppi} \sin(nx) \cot \frac{x}{2} \d x.
    \end{equation}
    记 $I_n = \int_{0}^{\uppi} \sin(nx) \cot \frac{x}{2} \d x$．当 $n \geqslant 2$ 时，利用三角恒等式有
    \begin{equation}\label{eq:trig-identity}
        \sin(nx) - \sin((n-1)x) = 2 \cos\left( \left(n - \frac{1}{2}\right)x \right) \sin \frac{x}{2},
    \end{equation}
    从而得到
    \begin{equation}\label{eq:diff-identity}
        \sin(nx) \cot \frac{x}{2} - \sin((n-1)x) \cot \frac{x}{2} = 2 \cos\left( \left(n - \frac{1}{2}\right)x \right) \cos \frac{x}{2} = \cos(nx) + \cos((n-1)x).
    \end{equation}
    对其在 $[0, \uppi]$ 上积分，由于当 $k \geqslant 1$ 时 $\int_{0}^{\uppi} \cos(kx) \d x = 0$，故对 $n \geqslant 2$ 有
    \begin{equation}\label{eq:in-diff}
        I_n - I_{n-1} = \int_{0}^{\uppi} \left[ \cos(nx) + \cos((n-1)x) \right] \d x = 0.
    \end{equation}
    而对于 $n=1$：
    \begin{equation}\label{eq:i1-calc}
        I_1 = \int_{0}^{\uppi} \sin x \cot \frac{x}{2} \d x = \int_{0}^{\uppi} 2 \cos^2 \frac{x}{2} \d x = \int_{0}^{\uppi} (1 + \cos x) \d x = \uppi.
    \end{equation}
    因此对所有 $n \geqslant 1$，均有 $I_n = \uppi$．将其代入\zcref{eq:an-parts}，得
    \begin{equation}\label{eq:an-final}
        a_n = \frac{1}{n\uppi} \cdot \uppi = \frac{1}{n}.
    \end{equation}
    由于 $f(x)$ 在 $(0, 2\uppi)$ 内连续且分段光滑，由 Fourier 级数的收敛定理，在其连续点处级数收敛于函数值，即对 $0 < x < 2\uppi$，有
    \begin{equation}\label{eq:res-2-mid}
        -\ln \left( 2\sin \frac{x}{2} \right) = \sum_{n=1}^{\infty} \frac{\cos nx}{n}.
    \end{equation}
    即得
    \begin{equation}
        \boxed{\sum_{n=1}^{\infty} \frac{\cos nx}{n} = - \ln \left( 2\sin \frac{x}{2} \right)}
    \end{equation}

    \ansitem{1}

    在已证的\zcref{eq:res-2-mid}中，对 $-\uppi < x < \uppi$，令 $y = \uppi - x$．此时显然有 $0 < y < 2\uppi$．
    将 $y$ 代入\zcref{eq:res-2-mid}，得
    \begin{equation}\label{eq:sub-step}
        -\ln \left( 2\sin \frac{\uppi - x}{2} \right) = \sum_{n=1}^{\infty} \frac{\cos(n(\uppi - x))}{n}.
    \end{equation}
    利用三角诱导公式 $\sin\left(\frac{\uppi}{2} - \frac{x}{2}\right) = \cos \frac{x}{2}$，以及 $\cos(n\uppi - nx) = (-1)^n \cos(nx)$，上式可化为
    \begin{equation}\label{eq:sub-simplify}
        -\ln \left( 2\cos \frac{x}{2} \right) = \sum_{n=1}^{\infty} (-1)^n \frac{\cos nx}{n} = - \sum_{n=1}^{\infty} (-1)^{n-1} \frac{\cos nx}{n}.
    \end{equation}
    两边同乘以 $-1$，即得
    \begin{equation}
        \boxed{\sum_{n=1}^{\infty} (-1)^{n-1} \frac{\cos nx}{n} = \ln \left( 2\cos \frac{x}{2} \right)}
    \end{equation}
\end{myproof}

\begin{problem}{Fourier系数的符号}{Fourier-Coefficient-Signs}
    设 $f$ 是周期为 $2\uppi$ 的可积且绝对可积函数．证明：
    \begin{enumerate}
        \item 如果 $f$ 在 $(0, 2\uppi)$ 上递减，那么 $b_n \geqslant 0$；
        \item 如果 $f$ 在 $(0, 2\uppi)$ 上递增，那么 $b_n \leqslant 0$．
    \end{enumerate}
\end{problem}

\begin{myproof}
    \ansitem{1}

    因为 $f$ 的 Fourier 系数 $b_n$ 定义为
    \begin{equation}\label{eq:bn-def-3}
        b_n = \frac{1}{\uppi} \int_{0}^{2\uppi} f(x) \sin(nx) \d x,
    \end{equation}
    我们将积分区间 $[0, 2\uppi]$ 划分为 $2n$ 个长度为 $\frac{\uppi}{n}$ 的子区间，即有
    \begin{equation}\label{eq:bn-split-3}
        \begin{aligned}
            b_n & = \frac{1}{\uppi} \sum_{k=0}^{2n-1} \int_{\frac{k\uppi}{n}}^{\frac{(k+1)\uppi}{n}} f(x) \sin(nx) \d x                                                        \\
                & = \frac{1}{\uppi} \sum_{k=0}^{2n-1} \int_{0}^{\frac{\uppi}{n}} f\left( t + \frac{k\uppi}{n} \right) \sin\left(n\left(t + \frac{k\uppi}{n}\right)\right) \d t \\
                & = \frac{1}{\uppi} \sum_{k=0}^{2n-1} (-1)^k \int_{0}^{\frac{\uppi}{n}} f\left( t + \frac{k\uppi}{n} \right) \sin(nt) \d t.
        \end{aligned}
    \end{equation}
    将 \zcref{eq:bn-split-3} 中的奇数项与偶数项配对，可得
    \begin{equation}\label{eq:bn-paired-3}
        b_n = \frac{1}{\uppi} \int_{0}^{\frac{\uppi}{n}} \left[ \sum_{j=0}^{n-1} \left( f\left( t + \frac{2j\uppi}{n} \right) - f\left( t + \frac{(2j+1)\uppi}{n} \right) \right) \right] \sin(nt) \d t.
    \end{equation}
    对于任意的 $t \in \left(0, \frac{\uppi}{n}\right)$ 和 $j \in \{0, 1, \dots, n-1\}$，由于
    \begin{equation}\label{eq:interval-check-3}
        0 < t + \frac{2j\uppi}{n} < t + \frac{(2j+1)\uppi}{n} < \frac{\uppi}{n} + \frac{(2n-1)\uppi}{n} = 2\uppi,
    \end{equation}
    可知自变量均落在区间 $(0, 2\uppi)$ 内．因为 $f$ 在 $(0, 2\uppi)$ 上递减，所以有
    \begin{equation}
        f\left( t + \frac{2j\uppi}{n} \right) \geqslant f\left( t + \frac{(2j+1)\uppi}{n} \right).
    \end{equation}
    又在积分区间 $t \in \left(0, \frac{\uppi}{n}\right)$ 上有 $\sin(nt) > 0$，故被积函数非负．因此由 \zcref{eq:bn-paired-3} 可知 $b_n \geqslant 0$．

    \ansitem{2}

    同理，若 $f$ 在 $(0, 2\uppi)$ 上递增，则对任意的 $t \in \left(0, \frac{\uppi}{n}\right)$ 和 $j \in \{0, 1, \dots, n-1\}$，有
    \begin{equation}
        f\left( t + \frac{2j\uppi}{n} \right) \leqslant f\left( t + \frac{(2j+1)\uppi}{n} \right),
    \end{equation}
    从而有
    \begin{equation}
        f\left( t + \frac{2j\uppi}{n} \right) - f\left( t + \frac{(2j+1)\uppi}{n} \right) \leqslant 0.
    \end{equation}
    由于 $\sin(nt) > 0$，因此 \zcref{eq:bn-paired-3} 中被积函数非正，即得 $b_n \leqslant 0$．
\end{myproof}

\begin{problem}{Fourier系数的渐近估计}{Fourier-Coefficient-Asymptotics}
    设 $f$ 是周期为 $2\uppi$ 且在 $[-\uppi, \uppi]$ 上 Riemann 可积的函数．如果它在 $(-\uppi, \uppi)$ 上单调，证明：
    \begin{equation}
        a_n = O\left(\frac{1}{n}\right), \quad b_n = O\left(\frac{1}{n}\right) \quad (n \to \infty).
    \end{equation}
\end{problem}

\begin{myproof}
    由于 $f$ 在 $[-\uppi, \uppi]$ 上 Riemann 可积，故 $f$ 在该区间上有界，从而其单侧极限 $f(-\uppi^+)$ 与 $f(\uppi^-)$ 均存在且有限．

    不妨设 $f$ 在 $(-\uppi, \uppi)$ 上单调递增．由积分第二中值定理，存在 $\xi_n \in [-\uppi, \uppi]$，使得
    \begin{equation}\label{eq:an-integral-4}
        \begin{aligned}
            \int_{-\uppi}^{\uppi} f(x) \cos(nx) \d x & = f(-\uppi^+) \int_{-\uppi}^{\xi_n} \cos(nx) \d x + f(\uppi^-) \int_{\xi_n}^{\uppi} \cos(nx) \d x                                     \\
                                                     & = f(-\uppi^+) \left[ \frac{\sin(n\xi_n) - \sin(-n\uppi)}{n} \right] + f(\uppi^-) \left[ \frac{\sin(n\uppi) - \sin(n\xi_n)}{n} \right] \\
                                                     & = \left[ f(-\uppi^+) - f(\uppi^-) \right] \frac{\sin(n\xi_n)}{n}.
        \end{aligned}
    \end{equation}
    因为 $f(-\uppi^+)$ 与 $f(\uppi^-)$ 为有限常数，且 $|\sin(n\xi_n)| \leqslant 1$，由 \zcref{eq:an-integral-4} 可得
    \begin{equation}
        |a_n| = \left| \frac{1}{\uppi} \int_{-\uppi}^{\uppi} f(x) \cos(nx) \d x \right| \leqslant \frac{|f(-\uppi^+) - f(\uppi^-)|}{\uppi n}.
    \end{equation}
    因此，当 $n \to \infty$ 时，有
    \begin{equation}
        a_n = O\left(\frac{1}{n}\right).
    \end{equation}

    同理，由积分第二中值定理，存在 $\eta_n \in [-\uppi, \uppi]$，使得
    \begin{equation}\label{eq:bn-integral-4}
        \begin{aligned}
            \int_{-\uppi}^{\uppi} f(x) \sin(nx) \d x & = f(-\uppi^+) \int_{-\uppi}^{\eta_n} \sin(nx) \d x + f(\uppi^-) \int_{\eta_n}^{\uppi} \sin(nx) \d x                        \\
                                                     & = f(-\uppi^+) \left[ \frac{(-1)^n - \cos(n\eta_n)}{n} \right] + f(\uppi^-) \left[ \frac{\cos(n\eta_n) - (-1)^n}{n} \right] \\
                                                     & = \left[ f(-\uppi^+) - f(\uppi^-) \right] \frac{(-1)^n - \cos(n\eta_n)}{n}.
        \end{aligned}
    \end{equation}
    因为 $|(-1)^n - \cos(n\eta_n)| \leqslant 2$，由 \zcref{eq:bn-integral-4} 可得
    \begin{equation}
        |b_n| = \left| \frac{1}{\uppi} \int_{-\uppi}^{\uppi} f(x) \sin(nx) \d x \right| \leqslant \frac{2|f(-\uppi^+) - f(\uppi^-)|}{\uppi n}.
    \end{equation}
    因此，当 $n \to \infty$ 时，有
    \begin{equation}
        b_n = O\left(\frac{1}{n}\right).
    \end{equation}
    若 $f$ 在 $(-\uppi, \uppi)$ 上单调递减，同理可证．
\end{myproof}

\begin{problem}{含参量积分的极限}{Limit-of-Parametric-Integral}
    设 $f$ 在 $[-a, a]$ 上连续，且在 $x = 0$ 处可导．求证：
    \begin{equation}
        \lim_{n \to +\infty} \int_{-a}^{a} \frac{1 - \cos n x}{x} f(x) \d x = \int_{0}^{a} \frac{f(x) - f(-x)}{x} \d x.
    \end{equation}
\end{problem}

\begin{myproof}
    令 $I(n) = \int_{-a}^{a} \frac{1 - \cos n x}{x} f(x) \d x$．对积分区间进行对称拆分，可得
    \begin{equation}\label{eq:split-5}
        \begin{aligned}
            I(n) & = \int_{-a}^{0} \frac{1 - \cos n x}{x} f(x) \d x + \int_{0}^{a} \frac{1 - \cos n x}{x} f(x) \d x      \\
                 & = \int_{0}^{a} \frac{1 - \cos n(-t)}{-t} f(-t) \d(-t) + \int_{0}^{a} \frac{1 - \cos n x}{x} f(x) \d x \\
                 & = -\int_{0}^{a} \frac{1 - \cos n t}{t} f(-t) \d t + \int_{0}^{a} \frac{1 - \cos n x}{x} f(x) \d x     \\
                 & = \int_{0}^{a} \frac{f(x) - f(-x)}{x} (1 - \cos n x) \d x.
        \end{aligned}
    \end{equation}
    定义辅助函数 $g(x)$ 如下：
    \begin{equation}
        g(x) = \begin{dcases}
            \frac{f(x) - f(-x)}{x}, & x \in (0, a], \\
            2f'(0),                 & x = 0.
        \end{dcases}
    \end{equation}
    由于 $f$ 在 $[-a, a]$ 上连续且在 $x = 0$ 处可导，有
    \begin{equation}
        \lim_{x \to 0^+} g(x) = \lim_{x \to 0^+} \frac{f(x) - f(0) - [f(-x) - f(0)]}{x} = f'(0) - [-f'(0)] = 2f'(0) = g(0).
    \end{equation}
    因此 $g(x)$ 在 $[0, a]$ 上连续，从而 Riemann 可积．由 \zcref{eq:split-5}，我们有
    \begin{equation}
        I(n) = \int_{0}^{a} g(x) \d x - \int_{0}^{a} g(x) \cos n x \d x.
    \end{equation}
    根据 Riemann-Lebesgue 引理，有
    \begin{equation}\label{eq:riemann-lebesgue-5}
        \lim_{n \to +\infty} \int_{0}^{a} g(x) \cos n x \d x = 0.
    \end{equation}
    结合 \zcref{eq:riemann-lebesgue-5}，取极限可得
    \begin{equation}
        \lim_{n \to +\infty} I(n) = \int_{0}^{a} g(x) \d x = \int_{0}^{a} \frac{f(x) - f(-x)}{x} \d x.
    \end{equation}
\end{myproof}

\begin{problem}{Fejer引理的推广}{Fejer-Lemma-Extension}
    设 $f$ 在 $[a, b]$ 上 Riemann 可积．证明：
    \begin{equation}
        \lim_{\lambda \to +\infty} \int_{a}^{b} f(x)|\cos \lambda x| \d x = \frac{2}{\uppi} \int_{a}^{b} f(x) \d x,
    \end{equation}
    \begin{equation}
        \lim_{\lambda \to +\infty} \int_{a}^{b} f(x)|\sin \lambda x| \d x = \frac{2}{\uppi} \int_{a}^{b} f(x) \d x.
    \end{equation}
\end{problem}

\begin{myproof}
    由于 $|\cos t|$ 和 $|\sin t|$ 均是以 $\uppi$ 为周期的 Riemann 可积函数，其在一个周期内的平均值均为
    \begin{equation}\label{eq:avg-6}
        M = \frac{1}{\uppi} \int_{0}^{\uppi} |\cos t| \d t = \frac{1}{\uppi} \int_{0}^{\uppi} |\sin t| \d t = \frac{2}{\uppi}.
    \end{equation}
    设 $g(t)$ 是以 $T$ 为周期的 Riemann 可积函数，且其平均值为 $M = \frac{1}{T} \int_{0}^{T} g(t) \d t$．我们先证明对任意 $[a, b]$ 上的 Riemann 可积函数 $f(x)$，有
    \begin{equation}\label{eq:fejer-general-6}
        \lim_{\lambda \to +\infty} \int_{a}^{b} f(x)g(\lambda x) \d x = M \int_{a}^{b} f(x) \d x.
    \end{equation}
    首先，若 $f(x) = 1$，则
    \begin{equation}
        \int_{a}^{b} g(\lambda x) \d x = \frac{1}{\lambda} \int_{\lambda a}^{\lambda b} g(u) \d u.
    \end{equation}
    设 $\lambda b - \lambda a = n T + \delta$，其中 $n = \lfloor \frac{\lambda(b-a)}{T} \rfloor \in \symbb{Z}$， $0 \leqslant \delta < T$．利用 $g$ 的周期性，可得
    \begin{equation}
        \begin{aligned}
            \int_{\lambda a}^{\lambda b} g(u) \d u & = n \int_{0}^{T} g(u) \d u + \int_{\lambda a}^{\lambda a + \delta} g(u) \d u \\
                                                   & = n T M + O(1).
        \end{aligned}
    \end{equation}
    因此
    \begin{equation}
        \lim_{\lambda \to +\infty} \int_{a}^{b} g(\lambda x) \d x = \lim_{\lambda \to +\infty} \left[ \frac{n T}{\lambda} M + O\left(\frac{1}{\lambda}\right) \right] = (b-a) M = M \int_{a}^{b} 1 \d x.
    \end{equation}
    由线性性质，\zcref{eq:fejer-general-6} 对任意阶梯函数均成立．

    对于一般的 Riemann 可积函数 $f(x)$，对任意 $\varepsilon > 0$，存在阶梯函数 $\varphi(x)$ 和 $\psi(x)$ 满足 $\varphi(x) \leqslant f(x) \leqslant \psi(x)$，且
    \begin{equation}
        \int_{a}^{b} (\psi(x) - \varphi(x)) \d x < \varepsilon.
    \end{equation}
    由于 $g(t) \geqslant 0$，我们有
    \begin{equation}
        \int_{a}^{b} \varphi(x) g(\lambda x) \d x \leqslant \int_{a}^{b} f(x) g(\lambda x) \d x \leqslant \int_{a}^{b} \psi(x) g(\lambda x) \d x.
    \end{equation}
    令 $\lambda \to +\infty$，由阶梯函数的结论可知
    \begin{equation}
        M \int_{a}^{b} \varphi(x) \d x \leqslant \liminf_{\lambda \to +\infty} \int_{a}^{b} f(x) g(\lambda x) \d x \leqslant \limsup_{\lambda \to +\infty} \int_{a}^{b} f(x) g(\lambda x) \d x \leqslant M \int_{a}^{b} \psi(x) \d x.
    \end{equation}
    又因为
    \begin{equation}
        M \int_{a}^{b} f(x) \d x - M \varepsilon < M \int_{a}^{b} \varphi(x) \d x,
    \end{equation}
    \begin{equation}
        M \int_{a}^{b} \psi(x) \d x < M \int_{a}^{b} f(x) \d x + M \varepsilon,
    \end{equation}
    由 $\varepsilon$ 的任意性，即得 \zcref{eq:fejer-general-6} 成立．

    将 $g(t) = |\cos t|$ 和 $g(t) = |\sin t|$ 分别代入 \zcref{eq:fejer-general-6}，并结合 \zcref{eq:avg-6}，即证
    \begin{equation}
        \boxed{\lim_{\lambda \to +\infty} \int_{a}^{b} f(x)| \cos \lambda x | \d x = \frac{2}{\uppi} \int_{a}^{b} f(x) \d x},
    \end{equation}
    \begin{equation}
        \boxed{\lim_{\lambda \to +\infty} \int_{a}^{b} f(x)| \sin \lambda x | \d x = \frac{2}{\uppi} \int_{a}^{b} f(x) \d x}.
    \end{equation}
\end{myproof}

\begin{problem}{卷积与Parseval等式}{Convolution-and-Parseval}
    设 $f$ 是周期为 $2\uppi$ 的连续函数．令
    \begin{equation}
        F(x) = \frac{1}{\uppi} \int_{-\uppi}^{\uppi} f(t)f(x + t) \d t,
    \end{equation}
    用 $a_n$、$b_n$ 和 $A_n$、$B_n$ 分别表示 $f$ 和 $F$ 的 Fourier 系数．证明：
    \begin{equation}
        A_0 = a_0^2, \quad A_n = a_n^2 + b_n^2, \quad B_n = 0.
    \end{equation}
    由此推出 $f$ 的 Parseval 等式．
\end{problem}

\begin{myproof}
    首先计算 $F(x)$ 的 Fourier 系数 $A_0$：
    \begin{equation}\label{eq:A0-calc-7}
        \begin{aligned}
            A_0 & = \frac{1}{\uppi} \int_{-\uppi}^{\uppi} F(x) \d x                                      \\
                & = \frac{1}{\uppi^2} \int_{-\uppi}^{\uppi} \int_{-\uppi}^{\uppi} f(t) f(x+t) \d t \d x.
        \end{aligned}
    \end{equation}
    交换积分顺序并作变量代换 $u = x+t$，由于 $f$ 的周期为 $2\uppi$，可得
    \begin{equation}
        \begin{aligned}
            A_0 & = \frac{1}{\uppi^2} \int_{-\uppi}^{\uppi} f(t) \left[ \int_{-\uppi}^{\uppi} f(u) \d u \right] \d t                              \\
                & = \left( \frac{1}{\uppi} \int_{-\uppi}^{\uppi} f(t) \d t \right) \left( \frac{1}{\uppi} \int_{-\uppi}^{\uppi} f(u) \d u \right) \\
                & = a_0^2.
        \end{aligned}
    \end{equation}

    接着计算 $A_n$：
    \begin{equation}\label{eq:An-calc-7}
        \begin{aligned}
            A_n & = \frac{1}{\uppi} \int_{-\uppi}^{\uppi} F(x) \cos(nx) \d x                                      \\
                & = \frac{1}{\uppi^2} \int_{-\uppi}^{\uppi} \int_{-\uppi}^{\uppi} f(t) f(x+t) \cos(nx) \d t \d x.
        \end{aligned}
    \end{equation}
    作变量代换 $u = x+t$，即 $x = u-t$，利用 $\cos n(u-t) = \cos(nu) \cos(nt) + \sin(nu) \sin(nt)$，得
    \begin{equation}
        \begin{aligned}
            A_n & = \frac{1}{\uppi^2} \int_{-\uppi}^{\uppi} \int_{-\uppi}^{\uppi} f(t) f(u) \left[ \cos(nu) \cos(nt) + \sin(nu) \sin(nt) \right] \d u \d t                \\
                & = \left( \frac{1}{\uppi} \int_{-\uppi}^{\uppi} f(u) \cos(nu) \d u \right) \left( \frac{1}{\uppi} \int_{-\uppi}^{\uppi} f(t) \cos(nt) \d t \right)       \\
                & \quad + \left( \frac{1}{\uppi} \int_{-\uppi}^{\uppi} f(u) \sin(nu) \d u \right) \left( \frac{1}{\uppi} \int_{-\uppi}^{\uppi} f(t) \sin(nt) \d t \right) \\
                & = a_n^2 + b_n^2.
        \end{aligned}
    \end{equation}

    再计算 $B_n$：
    \begin{equation}\label{eq:Bn-calc-7}
        \begin{aligned}
            B_n & = \frac{1}{\uppi} \int_{-\uppi}^{\uppi} F(x) \sin(nx) \d x                                       \\
                & = \frac{1}{\uppi^2} \int_{-\uppi}^{\uppi} \int_{-\uppi}^{\uppi} f(t) f(u) \sin n(u-t) \d u \d t.
        \end{aligned}
    \end{equation}
    利用 $\sin n(u-t) = \sin(nu) \cos(nt) - \cos(nu) \sin(nt)$，得
    \begin{equation}
        \begin{aligned}
            B_n & = \left( \frac{1}{\uppi} \int_{-\uppi}^{\uppi} f(u) \sin(nu) \d u \right) \left( \frac{1}{\uppi} \int_{-\uppi}^{\uppi} f(t) \cos(nt) \d t \right)       \\
                & \quad - \left( \frac{1}{\uppi} \int_{-\uppi}^{\uppi} f(u) \cos(nu) \d u \right) \left( \frac{1}{\uppi} \int_{-\uppi}^{\uppi} f(t) \sin(nt) \d t \right) \\
                & = b_n a_n - a_n b_n                                                                                                                                     \\
                & = 0.
        \end{aligned}
    \end{equation}
    由上述计算，即得
    \begin{equation}
        \boxed{A_0 = a_0^2, \quad A_n = a_n^2 + b_n^2, \quad B_n = 0}.
    \end{equation}

    下面推导 $f$ 的 Parseval 等式．
    由于 $f$ 在 $[-\uppi, \uppi]$ 上连续且周期为 $2\uppi$，作为连续函数的自我卷积， $F(x)$ 也是周期为 $2\uppi$ 的连续函数．由于 $B_n = 0$ 且其 Fourier 系数满足 $A_n = a_n^2 + b_n^2 \geqslant 0$，且 $\sum_{n=1}^{\infty} (a_n^2 + b_n^2) < \infty$，可知 $F(x)$ 的 Fourier 级数绝对且一致收敛于自身，即对任意 $x \in \symbb{R}$ 有
    \begin{equation}
        F(x) = \frac{A_0}{2} + \sum_{n=1}^{\infty} A_n \cos(nx).
    \end{equation}
    代入 $x = 0$，得
    \begin{equation}
        F(0) = \frac{A_0}{2} + \sum_{n=1}^{\infty} A_n.
    \end{equation}
    将 $A_0$ 与 $A_n$ 的表达式代入上式，并根据 $F(0)$ 的定义，即可推出 $f$ 的 Parseval 等式：
    \begin{equation}
        \boxed{\frac{1}{\uppi} \int_{-\uppi}^{\uppi} f^2(t) \d t = \frac{a_0^2}{2} + \sum_{n=1}^{\infty} (a_n^2 + b_n^2)}.
    \end{equation}
\end{myproof}

\begin{problem}{Wirtinger不等式}{Wirtinger-Inequality}
    设 $f$ 在 $[-\uppi, \uppi]$ 上连续，且在此区间上有可积且平方可积的导数 $f'$．如果 $f$ 满足
    \begin{equation}
        f(-\uppi) = f(\uppi), \quad \int_{-\uppi}^{\uppi} f(x) \d x = 0,
    \end{equation}
    证明：
    \begin{equation}
        \int_{-\uppi}^{\uppi} (f'(x))^2 \d x \geqslant \int_{-\uppi}^{\uppi} f^2(x) \d x,
    \end{equation}
    等号当且仅当 $f(x) = \alpha \cos x + \beta \sin x$ 时成立．
\end{problem}

\begin{myproof}
    由于 $f$ 在 $[-\uppi, \uppi]$ 上连续且 $f(-\uppi) = f(\uppi)$，可将其周期延拓为以 $2\uppi$ 为周期的连续函数．设其 Fourier 级数为
    \begin{equation}\label{eq:f-fourier-8}
        f(x) \sim \frac{a_0}{2} + \sum_{n=1}^{\infty} (a_n \cos(nx) + b_n \sin(nx)).
    \end{equation}
    根据已知条件 $\int_{-\uppi}^{\uppi} f(x) \d x = 0$，可知 $a_0 = 0$．因此，由 Parseval 等式有
    \begin{equation}\label{eq:f-parseval-8}
        \frac{1}{\uppi} \int_{-\uppi}^{\uppi} f^2(x) \d x = \sum_{n=1}^{\infty} (a_n^2 + b_n^2).
    \end{equation}

    由于 $f'$ 在 $[-\uppi, \uppi]$ 上可积且平方可积，且 $f(-\uppi) = f(\uppi)$，可以通过逐项求导得到 $f'$ 的 Fourier 展开式：
    \begin{equation}
        f'(x) \sim \sum_{n=1}^{\infty} (n b_n \cos(nx) - n a_n \sin(nx)).
    \end{equation}
    对 $f'$ 应用 Parseval 等式，可得
    \begin{equation}\label{eq:fprime-parseval-8}
        \frac{1}{\uppi} \int_{-\uppi}^{\uppi} (f'(x))^2 \d x = \sum_{n=1}^{\infty} n^2 (a_n^2 + b_n^2).
    \end{equation}
    比较 \zcref{eq:f-parseval-8} 与 \zcref{eq:fprime-parseval-8}，由于对所有 $n \geqslant 1$ 均有 $n^2 \geqslant 1$，可知
    \begin{equation}
        \sum_{n=1}^{\infty} n^2 (a_n^2 + b_n^2) \geqslant \sum_{n=1}^{\infty} (a_n^2 + b_n^2),
    \end{equation}
    即
    \begin{equation}
        \frac{1}{\uppi} \int_{-\uppi}^{\uppi} (f'(x))^2 \d x \geqslant \frac{1}{\uppi} \int_{-\uppi}^{\uppi} f^2(x) \d x.
    \end{equation}
    由此即得
    \begin{equation}
        \boxed{\int_{-\uppi}^{\uppi} (f'(x))^2 \d x \geqslant \int_{-\uppi}^{\uppi} f^2(x) \d x}.
    \end{equation}

    等号成立当且仅当对所有 $n \geqslant 1$，有
    \begin{equation}
        n^2 (a_n^2 + b_n^2) = a_n^2 + b_n^2.
    \end{equation}
    对于 $n \geqslant 2$，由于 $n^2 > 1$，上式成立当且仅当 $a_n = b_n = 0$．
    因此，等号成立的充要条件是仅有 $n=1$ 处的系数可能不为零，此时 $f(x)$ 具有形式
    \begin{equation}
        f(x) = a_1 \cos x + b_1 \sin x.
    \end{equation}
    令 $\alpha = a_1$， $\beta = b_1$，即等号当且仅当 $f(x) = \alpha \cos x + \beta \sin x$ 时成立．
\end{myproof}

\endinput
